% Thesis insertion: connection between the low-rank toy and LLM phase dynamics
\subsection{Why a low-rank toy is informative for LLM pre-training dynamics}
\label{sec:toy_llm_connection}

The low-rank toy is not intended to be an architectural miniature of a transformer. Its role is instead to isolate a candidate dynamical mechanism that is difficult to observe directly in a pretrained language model: the formation, compression, and consolidation of data-aligned latent feature subspaces. The relevant analogy is therefore not architectural homology, but dynamical homology. In the Pythia experiments, the objects of study are weight matrices inside a transformer trained by next-token prediction; in the toy, the object of study is a bilinear embedding model trained on observations from a known low-rank teacher. These systems differ in surface form, but they share the same abstract problem: a learner must infer a lower-dimensional latent structure from many partial observations, and its internal weights should change from random high-dimensional structure toward a smaller number of data-aligned directions.

The toy task makes this abstraction explicit. A teacher samples latent entity features $U_e$ and latent relation features $V_r$, then generates observations
\[
    y_{e,r} = U_e^\top V_r + \epsilon.
\]
The student receives only a subset of the entries $(e,r,y_{e,r})$ and learns embeddings $\hat U_e$ and $\hat V_r$ by minimizing prediction error. Held-out matrix-completion performance therefore depends on recovering the teacher's latent subspace rather than memorising individual observations. This is the property that the earlier discrete toy tasks lacked: in those settings the model could fit training examples while held-out generalisation remained close to chance, leaving no reliable ground-truth transition to calibrate the spectral indicators.

The connection to language modelling is that next-token prediction also contains many partial observations of latent structure. Contexts, tokens, facts, syntactic patterns, topics, and behaviours are not learned as isolated strings; they are tied together by repeated statistical regularities. A transformer is far more expressive than the bilinear toy, but if part of its early training consists of discovering reusable feature directions from such regularities, then one expects a qualitative sequence similar to the toy: initially unstable directions, emergence of data-aligned outliers, compression into lower-rank dominant structure, and later stabilisation of those directions. E1 measures this sequence in Pythia weight space. The toy supplies a setting where the true latent subspace is known, so that the same spectral signatures can be compared against an independent ground truth.

This distinction also fixes the interpretation of the toy. A positive toy result does not prove that Pythia literally performs matrix completion, nor that a transformer implements the same bilinear computation. It shows a sufficiency claim: when a learner recovers and consolidates a latent feature subspace, the same family of indicators used in E1 can exhibit phase-like dynamics, and intervention durability can change across the transition. The toy therefore justifies the plausibility of the mechanism used to interpret the Pythia results, while the Pythia experiments remain the empirical evidence for LLM pre-training dynamics.

\paragraph{Mapping between the toy and the Pythia experiments.}
The conceptual mapping is as follows. The teacher latent features in the toy correspond to reusable statistical features in language data. Observed matrix entries correspond to partial observations of these features in training examples. Held-out matrix-completion error corresponds to functional generalisation beyond memorised samples. Stable-rank collapse and subspace alignment in the toy correspond to the spectral compression and singular-subspace stabilisation measured in E1. Frequency imbalance in the toy corresponds to the fact that some linguistic or factual regularities occur with higher signal-to-noise than others. Injected exception entries correspond to the synthetic factual associations used in E3. Clean continuation corresponds to further training on the original data distribution, while conflicting exceptions correspond to the adversarial overwrite/degradation branch.

The toy is therefore used as a controlled mechanism vignette. It strengthens the thesis only if it satisfies three conditions: held-out performance improves, the learned subspace aligns with the known teacher subspace, and the spectral indicators detect the same transition. Under those conditions, the toy provides an interpretable explanation for why Pythia might show an early reorganisation window: the model is not merely increasing weight norm, but forming and consolidating data-aligned feature directions.
